% !TeX program = pdflatex
\documentclass[10pt, aspectratio=169]{beamer}

% Preámbulo modular para presentaciones Beamer (Modelación Estadística)
% Diseño y paleta institucional del Tecnológico de Monterrey

\documentclass[aspectratio=169,xcolor={dvipsnames,table}]{beamer}

\usepackage[utf8]{inputenc}
\usepackage[spanish,mexico]{babel}
\usepackage{amsmath,amssymb,amsfonts,mathtools}
\usepackage{graphicx}
\usepackage{listings}
\usepackage{booktabs}
\usepackage{tikz}
\usetikzlibrary{matrix,arrows,decorations.pathmorphing,shapes.geometric,positioning}

% Paleta Institucional Tecnológico de Monterrey
\definecolor{TecAzul}{HTML}{0039A6}       % Azul TEC: color primario institucional
\definecolor{TecAzulOscuro}{HTML}{1A2E51} % Azul marino oscuro
\definecolor{TecRojo}{HTML}{EC2661}       % Rojo de acento (paleta miTec)
\definecolor{TecGris}{HTML}{646464}       % Gris neutro para comentarios y subtítulos
\definecolor{TecFondoBloque}{HTML}{F4F6F9}% Fondo suave para bloques

% Configuración de Tema Beamer
\usetheme{Boadilla}
\usecolortheme[named=TecAzulOscuro]{structure}
\setbeamercolor{alerted text}{fg=TecRojo}
\setbeamercolor{palette primary}{bg=TecAzulOscuro,fg=white}
\setbeamercolor{palette secondary}{bg=TecAzul,fg=white}
\setbeamercolor{palette tertiary}{bg=TecAzulOscuro!80!black,fg=white}
\setbeamercolor{title}{fg=TecAzulOscuro,bg=TecFondoBloque}
\setbeamercolor{frametitle}{fg=TecAzulOscuro,bg=TecFondoBloque}

% Bloques personalizados elegantes
\setbeamercolor{block title}{bg=TecAzulOscuro,fg=white}
\setbeamercolor{block body}{bg=TecFondoBloque,fg=black}
\setbeamercolor{block title alerted}{bg=TecRojo,fg=white}
\setbeamercolor{block body alerted}{bg=TecRojo!8,fg=black}
\setbeamercolor{block title example}{bg=TecAzul,fg=white}
\setbeamercolor{block body example}{bg=TecAzul!8,fg=black}

% Redondear bordes y añadir sombra sutil
\setbeamertemplate{blocks}[rounded][shadow=true]
\setbeamertemplate{navigation symbols}{}

% Pie de página limpio con número de diapositiva
\setbeamertemplate{footline}{
  \leavevmode%
  \hbox{%
  \begin{beamercolorbox}[wd=.7\paperwidth,ht=2.25ex,dp=1ex,left]{author in head/foot}%
    \hspace*{2ex}\insertshortauthor~~(\insertshortinstitute) \hfill \insertshorttitle
  \end{beamercolorbox}%
  \begin{beamercolorbox}[wd=.3\paperwidth,ht=2.25ex,dp=1ex,right]{date in head/foot}%
    \insertshortdate{}\hspace*{2em}
    \insertframenumber{} / \inserttotalframenumber\hspace*{2ex}
  \end{beamercolorbox}}%
  \vskip0pt%
}

% Configuración de Listings de Python para visualización pedagógica de código
\lstset{
    language=Python,
    basicstyle=\ttfamily\footnotesize,
    keywordstyle=\color{TecAzulOscuro}\bfseries,
    stringstyle=\color{TecRojo},
    commentstyle=\color{TecGris}\itshape,
    showstringspaces=false,
    breaklines=true,
    frame=lines,
    rulecolor=\color{TecAzulOscuro!30},
    backgroundcolor=\color{TecFondoBloque},
    aboveskip=0.5em,
    belowskip=0.5em,
    literate={á}{{\'a}}1 {é}{{\'e}}1 {í}{{\'i}}1 {ó}{{\'o}}1 {ú}{{\'u}}1
             {Á}{{\'A}}1 {É}{{\'E}}1 {Í}{{\'I}}1 {Ó}{{\'O}}1 {Ú}{{\'U}}1
             {ñ}{{\~n}}1 {Ñ}{{\~N}}1 {¿}{{?`}}1 {¡}{{!`}}1
}

% Metadatos institucionales por defecto
\title[Modelación Estadística]{Modelación Estadística para Ciencia de Datos e Ingeniería}
\author[J. Castillo Colmenares]{Juliho David Castillo Colmenares}
\institute[Tec de Monterrey]{Tecnológico de Monterrey \\ Escuela de Ingeniería y Ciencias}
\date{\today}

% Comandos y macros estadísticas compartidas para las presentaciones Beamer (Metropolis)

% Conjuntos numéricos básicos
\newcommand{\R}{\mathbb{R}}
\newcommand{\N}{\mathbb{N}}
\newcommand{\Q}{\mathbb{Q}}
\newcommand{\Z}{\mathbb{Z}}
\newcommand{\set}[1]{\left\{ #1 \right\}}
\newcommand{\abs}[1]{\left|#1\right|}
\newcommand{\norm}[1]{\left\| #1 \right\|}

% Comandos estadísticos y probabilísticos del curso
\newcommand{\Var}{\operatorname{Var}}
\newcommand{\cov}{\operatorname{Cov}}
\newcommand{\corr}{\rho}
\newcommand{\comb}[2]{\begin{pmatrix} #1 \\ #2 \end{pmatrix}}
\newcommand{\s}{\sigma}
\newcommand{\particion}{\mathfrak{P}}
\newcommand{\card}[1]{\#\left( #1 \right)}
\newcommand{\seq}[3]{\set{#1}_{#2}^{#3}}
\newcommand{\E}{\mathbb{E}}
\newcommand{\Prob}{\mathbb{P}}

% Entornos pedagógicos y cajas en español académico natural
\newenvironment{discusion}{%
  \begin{alertblock}{Pregunta de Discusión}%
}{%
  \end{alertblock}%
}

\newenvironment{reflexion}{%
  \begin{alertblock}{Reflexión para el Aula}%
}{%
  \end{alertblock}%
}

\newenvironment{paradoja}{%
  \begin{alertblock}{Paradoja de Exploración}%
}{%
  \end{alertblock}%
}

\newenvironment{motivacion}{%
  \begin{exampleblock}{Motivación: Ciencia de Datos en la Práctica}%
}{%
  \end{exampleblock}%
}

\newenvironment{retoclase}{%
  \begin{block}{Reto Rápido en Equipo}%
}{%
  \end{block}%
}


\title[Bernoulli y Binomial]{Sección 03.04: Bernoulli y Binomial}
\subtitle{Procesos Dicotómicos e Inferencia Combinatoria}
\author[J. Castillo Colmenares]{Juliho Castillo Colmenares}
\institute{Tecnológico de Monterrey}
\date{\vspace{-2.0cm}}

\begin{document}

% Slide 1: Portada
\begin{frame}[plain]
  \vspace{-0.5cm}
  \titlepage
\end{frame}

% Slide 2: Hoja de Ruta
\begin{frame}{Hoja de Ruta — Capítulo 03: Variables Aleatorias Discretas}
  \begin{block}{Estructura Curricular del Capítulo 03}
    \footnotesize
    \begin{enumerate}\itemsep=0.08cm
      \item \textcolor{gray}{Sección 03.01: Funciones de Probabilidad Discretas (PMF y Soporte)}
      \item \textcolor{gray}{Sección 03.02: Función de Distribución Acumulada (CDF Discreta)}
      \item \textcolor{gray}{Sección 03.03: Esperanza Matemática, Varianza y Momentos}
      \item \textbf{\textcolor{TecRojo}{Sección 03.04: Distribuciones de Bernoulli y Binomial}} $\leftarrow$ \emph{Foco actual}
      \item \textcolor{gray}{Sección 03.05: Distribuciones Geométrica y Binomial Negativa}
      \item \textcolor{gray}{Sección 03.06: Distribución Hipergeométrica}
    \end{enumerate}
  \end{block}
  \vspace{-0.05cm}
  \pause
  \begin{alertblock}{Objetivo Curricular de la Sección}
    \footnotesize
    Dominar la fundamentación algebraica, propiedades de momentos y modelación combinatoria de ensayos dicotómicos independientes, aplicando la distribución Binomial en control de calidad e ingeniería de sistemas.
  \end{alertblock}
\end{frame}

% Slide 3: Motivación e Intuición
\begin{frame}{Motivación — De la Dicotomía Elemental al Muestreo Repetido}
  \begin{columns}[T]
    \begin{column}{0.48\textwidth}
      \begin{block}{El Ladrillo Elemental de la Probabilidad}
        \footnotesize
        \begin{itemize}\itemsep=0.1cm
          \item En ciencia y tecnología, innumerables fenómenos se reducen a dos estados incompatibles: \textbf{Éxito / Fracaso}, \textbf{Apto / Defectuoso}, \textbf{1 / 0}.
          \item Un experimento elemental con dos salidas posibles se denomina \textbf{Ensayo de Bernoulli}.
          \item Modelar un solo ensayo permite cuantificar el riesgo elemental de un componente individual.
        \end{itemize}
      \end{block}
    \end{column}
    \begin{column}{0.48\textwidth}
      \pause
      \begin{alertblock}{La Necesidad de la Agregación}
        \footnotesize
        \begin{itemize}\itemsep=0.1cm
          \item La ingeniería rara vez evalúa piezas aisladas: inspecciona lotes de tamaño $n$, transmite paquetes de $n$ bits o realiza $n$ ensayos de estrés.
          \item Al repetir $n$ ensayos de Bernoulli de forma \textbf{independiente}, surge de manera natural la \textbf{Distribución Binomial}, que cuenta el número de éxitos acumulados en el sistema.
        \end{itemize}
      \end{alertblock}
    \end{column}
  \end{columns}
\end{frame}

% Slide 4: El Ensayo de Bernoulli
\begin{frame}{El Ensayo de Bernoulli — Definición Formal}
  \begin{block}{Variable Aleatoria de Bernoulli ($X \sim \text{Bernoulli}(p)$)}
    \footnotesize
    Una variable aleatoria discreta $X$ sigue una distribución de Bernoulli con parámetro de éxito $p \in (0, 1)$ si su soporte es $\mathcal{S}_X = \{0, 1\}$ y su función de masa de probabilidad está dada por:
    $$f_X(x) = P(X = x) = p^x (1-p)^{1-x}, \quad \text{para } x \in \{0, 1\}.$$
    Donde $q = 1-p$ representa la probabilidad complementaria de fracaso ($X=0$).
  \end{block}
  \vspace{-0.05cm}
  \pause
  \begin{columns}[T]
    \begin{column}{0.48\textwidth}
      \begin{exampleblock}{Propiedad Indicadora}
        \footnotesize
        $X$ actúa como la función indicadora de un evento $A \subset \Omega$: toma el valor $1$ si $\omega \in A$ y $0$ si $\omega \notin A$.
      \end{exampleblock}
    \end{column}
    \begin{column}{0.48\textwidth}
      \begin{alertblock}{Normalización Exacta}
        \footnotesize
        La suma sobre el soporte verifica la probabilidad total: $f_X(0) + f_X(1) = (1-p) + p = 1$.
      \end{alertblock}
    \end{column}
  \end{columns}
\end{frame}

% Slide 5: Momentos del Ensayo de Bernoulli
\begin{frame}{Momentos del Ensayo de Bernoulli}
  \begin{block}{Primer y Segundo Momento Bruto}
    \footnotesize
    Aplicando la definición de valor esperado al soporte dicotómico $\{0, 1\}$:
    $$\E[X] = \sum_{x=0}^1 x \cdot P(X=x) = 0 \cdot (1-p) + 1 \cdot p = p$$
    $$\E[X^2] = \sum_{x=0}^1 x^2 \cdot P(X=x) = 0^2 \cdot (1-p) + 1^2 \cdot p = p$$
  \end{block}
  \vspace{-0.05cm}
  \pause
  \begin{columns}[T]
    \begin{column}{0.48\textwidth}
      \begin{alertblock}{Varianza y Desviación Estándar}
        \footnotesize
        $$\Var(X) = \E[X^2] - (\E[X])^2 = p - p^2 = p(1-p) = pq$$
        $$\sigma_X = \sqrt{pq}$$
      \end{alertblock}
    \end{column}
    \begin{column}{0.48\textwidth}
      \begin{block}{Máxima Incertidumbre}
        \footnotesize
        La varianza alcanza su máximo absoluto en $p = 0.5$ ($\Var(X) = 0.25$). Si $p \to 0$ o $p \to 1$, el resultado se vuelve casi determinista y la varianza tiende a cero.
      \end{block}
    \end{column}
  \end{columns}
\end{frame}

% Slide 6: Génesis de la Distribución Binomial
\begin{frame}{Génesis de la Distribución Binomial — Suma de Ensayos Independientes}
  \begin{block}{Definición Constructiva ($X \sim \text{Bin}(n, p)$)}
    \footnotesize
    Sea $I_1, I_2, \dots, I_n$ una secuencia de $n$ ensayos de Bernoulli mutuamente independientes e idénticamente distribuidos (i.i.d.), donde $P(I_k = 1) = p$ para todo $k$. La variable aleatoria Binomial $X$ se define como la suma indicadora de éxitos:
    $$X = \sum_{k=1}^n I_k \implies \mathcal{S}_X = \{0, 1, 2, \dots, n\}.$$
  \end{block}
  \vspace{-0.05cm}
  \pause
  \begin{columns}[T]
    \begin{column}{0.48\textwidth}
      \begin{exampleblock}{Independencia y Probabilidad Conjunta}
        \footnotesize
        Una secuencia particular con $x$ éxitos y $n-x$ fracasos (ej. $1, 1, 0, 1, \dots, 0$) tiene probabilidad individual exacta $p^x (1-p)^{n-x}$ por independencia pura.
      \end{exampleblock}
    \end{column}
    \begin{column}{0.48\textwidth}
      \begin{alertblock}{El Factor Combinatorio}
        \footnotesize
        Existen exactamente $\comb{n}{x} = \dfrac{n!}{x!(n-x)!}$ formas distintas de permutar los $x$ éxitos en las $n$ posiciones del experimento.
      \end{alertblock}
    \end{column}
  \end{columns}
\end{frame}

% Slide 7: Función de Masa de Probabilidad Binomial
\begin{frame}{Función de Masa de Probabilidad Binomial (PMF)}
  \begin{block}{Teorema de la PMF Binomial}
    \footnotesize
    La función de masa de probabilidad de $X \sim \text{Bin}(n, p)$ es:
    $$f_X(x) = P(X = x) = \comb{n}{x} p^x (1-p)^{n-x}, \quad \text{para } x \in \{0, 1, \dots, n\}.$$
  \end{block}
  \vspace{-0.05cm}
  \pause
  \begin{columns}[T]
    \begin{column}{0.48\textwidth}
      \begin{exampleblock}{Normalización vía Teorema del Binomio}
        \footnotesize
        $$\sum_{x=0}^n \comb{n}{x} p^x q^{n-x} = (p + q)^n = 1^n = 1$$
        Esto garantiza formalmente la validez axiomática del modelo para todo $n \ge 1$ y $p \in (0, 1)$.
      \end{exampleblock}
    \end{column}
    \begin{column}{0.48\textwidth}
      \begin{alertblock}{Simetría y Complementariedad}
        \footnotesize
        Contar éxitos ($X$) con probabilidad $p$ es matemáticamente idéntico a contar fracasos ($Y = n - X$) con probabilidad $1-p$:
        $$X \sim \text{Bin}(n, p) \iff n-X \sim \text{Bin}(n, 1-p)$$
      \end{alertblock}
    \end{column}
  \end{columns}
\end{frame}

% Slide 8: Momentos y Asimetría de la Binomial
\begin{frame}{Momentos, Asimetría y Forma de la Distribución Binomial}
  \begin{columns}[T]
    \begin{column}{0.48\textwidth}
      \begin{block}{Aditividad de Momentos}
        \footnotesize
        Al ser $X = \sum_{k=1}^n I_k$ suma de independientes:
        $$\mu_X = \E[X] = \sum_{k=1}^n \E[I_k] = np$$
        $$\sigma_X^2 = \Var(X) = \sum_{k=1}^n \Var(I_k) = np(1-p)$$
        $$\sigma_X = \sqrt{np(1-p)}$$
      \end{block}
    \end{column}
    \begin{column}{0.48\textwidth}
      \pause
      \begin{alertblock}{Asimetría ($\gamma_1$) y Curtosis Excess ($\gamma_2$)}
        \footnotesize
        $$\gamma_1 = \dfrac{1 - 2p}{\sqrt{np(1-p)}}, \quad \gamma_2 = \dfrac{1 - 6p(1-p)}{np(1-p)}$$
        \begin{itemize}\itemsep=0.06cm
          \item Si $p = 0.5$, la PMF es perfectamente simétrica ($\gamma_1 = 0$).
          \item Si $p < 0.5$, asimetría positiva hacia la derecha; si $p > 0.5$, hacia la izquierda.
        \end{itemize}
      \end{alertblock}
    \end{column}
  \end{columns}
\end{frame}

% Slide 9A: Laboratorio en Python — PMF y CDF
\begin{frame}[fragile]{Laboratorio en Python: Validación Exacta PMF y CDF (`scipy.stats`)}
  \lstinputlisting[language=Python, firstline=12, lastline=36, basicstyle=\fontsize{6pt}{7pt}\selectfont\ttfamily]{../../code/03_variables_aleatorias_discretas/03.04_bernoulli_binomial.py}
\end{frame}

% Slide 9B: Laboratorio en Python — Monte Carlo
\begin{frame}[fragile]{Laboratorio en Python: Simulación Monte Carlo y Momentos}
  \lstinputlisting[language=Python, firstline=38, lastline=57, basicstyle=\fontsize{6.5pt}{7.5pt}\selectfont\ttfamily]{../../code/03_variables_aleatorias_discretas/03.04_bernoulli_binomial.py}
\end{frame}

% Slide 9C: Laboratorio en Python — Moda y Asimetría
\begin{frame}[fragile]{Laboratorio en Python: Verificación de Moda y Asimetría}
  \lstinputlisting[language=Python, firstline=59, lastline=75, basicstyle=\fontsize{6.5pt}{7.5pt}\selectfont\ttfamily]{../../code/03_variables_aleatorias_discretas/03.04_bernoulli_binomial.py}
\end{frame}

% Slide 9D: Laboratorio en Python — Salida en Terminal
\begin{frame}[fragile]{Laboratorio en Python: Salida en Terminal}
  \begin{block}{Ejecución en Consola (\texttt{python 03.04\_bernoulli\_binomial.py})}
    \fontsize{6pt}{7pt}\selectfont\ttfamily
=== Section 03.04: Bernoulli and Binomial Lab ===\\[0.08cm]
{[1]} Vectorized PMF \& CDF Validation (n=15, p=0.2):\\
\ \ \ \ Total Support Mass P(X<=n): 1.000000 (Target: 1.0)\\
\ \ \ \ Max PMF Error vs. Scipy\ \ : 6.66e-16\\
\ \ \ \ Max CDF Error vs. Scipy\ \ : 2.78e-15\\[0.08cm]
{[2]} Monte Carlo LLN Simulation (N=100,000):\\
\ \ \ \ Empirical Mean bar(X)\ \ \ \ : 2.9970 (Target: 3.0000, Err: 0.003020)\\
\ \ \ \ Empirical Variance S\^{}2\ \ \ \ : 2.3827 (Target: 2.4000, Err: 0.017289)\\[0.08cm]
{[3]} Mode Verification \& Skewness Analysis:\\
\ \ \ \ Computational Mode argmax : 3 (Theoretical floor((n+1)p): 3)\\
\ \ \ \ Exact Skewness gamma\_1\ \ \ \ : 0.387298 (Scipy Target: 0.387298)\\[0.08cm]
=== All numerical verifications passed exact tolerances. ===
  \end{block}
\end{frame}

% Slide 10A: Ejercicio en Clase — Nivel 1: Fundamental (Enunciado)

% Slide 10B: Ejercicio en Clase — Nivel 1: Fundamental (Resolución)

% Slide 11A: Ejercicio en Clase — Nivel 2: Operativo (Enunciado)

% Slide 11B: Ejercicio en Clase — Nivel 2: Operativo (Resolución)

% Slide 12A: Ejercicio en Clase — Nivel 3: Analítico (Enunciado)

% Slide 14: Cierre de Sección
\begin{frame}{Cierre de Sección y Síntesis Estocástica}
  \begin{block}{Resumen de Competencias en Modelación Binomial}
    \footnotesize
    \begin{itemize}\itemsep=0.1cm
      \item \textbf{Ensayo de Bernoulli ($X \sim \text{Bernoulli}(p)$):} El ladrillo elemental del análisis de riesgo e indicadores dicotómicos, con media $\mu = p$ y varianza máxima $\sigma^2 = pq \le 0.25$.
      \item \textbf{Agregación Combinatoria ($X \sim \text{Bin}(n, p)$):} El conteo de éxitos en $n$ ensayos independientes requiere el factor $\comb{n}{x}$ para dar cuenta de las permutaciones algebraicas del espacio muestral.
      \item \textbf{Aditividad de Momentos:} Por linealidad de sumas i.i.d., $\E[X] = np$ y $\Var(X) = npq$, sirviendo como patrón comparativo para detectar sobredispersión en modelos empíricos.
    \end{itemize}
  \end{block}
\end{frame}

% Slide 15: Perspectiva Modular
\begin{frame}{Perspectiva Modular — El Próximo Reto}
  \begin{columns}[T]
    \begin{column}{0.48\textwidth}
      \begin{block}{El Límite del Conteo Fijo ($n$ Constante)}
        \footnotesize
        Hasta ahora, el número de ensayos $n$ se ha mantenido \textbf{fijo de antemano}, mientras que el número de éxitos $X \in \{0, \dots, n\}$ es la variable aleatoria. ¿Qué ocurre si invertimos el paradigma experimental?
      \end{block}
    \end{column}
    \begin{column}{0.48\textwidth}
      \pause
      \begin{alertblock}{Hacia la Sección 03.05: Tiempos de Espera}
        \footnotesize
        En la \textbf{Sección 03.05}, fijaremos el número de éxitos deseados ($r \ge 1$) y dejaremos que el número de ensayos necesarios sea la variable aleatoria. Estudiaremos la \textbf{Distribución Geométrica} ($r=1$) y la \textbf{Binomial Negativa} ($r>1$), descubriendo la fascinante propiedad de pérdida de memoria.
      \end{alertblock}
    \end{column}
  \end{columns}
\end{frame}

\end{document}
